\section*{अध्याय \ref{ch:b2:affine} --- आनत समष्टियाँ}

\begin{solution}{exo:b2:affine:1}
$\{x + y + z = 1\}$: \omterm{def:b2:affine:subspace}{आनत} समतल, जिसकी दिशा सदिश समतल $\{x
+ y + z = 0\}$ और विमा $2$ है। $x - z = 3$ जोड़ने पर: कोई
\omterm{def:b2:affine:subspace}{आनत} रेखा (दो स्वतंत्र समीकरण), दिशा $\{x + y + z = 0,\ x = z\}
= \operatorname{Vect}\bigl((1, -2, 1)\bigr)$, विमा $1$। $\{x^2
+ y^2 = 1\}$: एक बेलन ---
बैरिकेंद्रों के अंतर्गत स्थायी नहीं ($(1,0,0)$ और $(-1,0,0)$ का मध्यबिंदु मूल बिंदु है,
जो बेलन पर नहीं): अतः \omterm{def:b2:affine:subspace}{आनत} नहीं। और कोई संगत निकाय $MX = B$: विमा $\dim\ker M$ की \omterm{def:b2:affine:subspace}{आनत}
उपसमष्टि $X_0 + \ker M$, जैसा \cref{def:b2:affine:subspace} में स्मरण किया गया।
\end{solution}

\begin{solution}{exo:b2:affine:2}
$C = \operatorname{bar}(A, 1-t; B, t)$ का अर्थ है $\vect{AC} =
t\,\vect{AB}$: और ऐसे $t$ का अस्तित्व ठीक $\vect{AC}$ का $\vect{AB} \neq 0$ के साथ जुड़ा
होना है, अर्थात् संरेखता। स्थितियाँ: $t
= \frac12$: मध्यबिंदु; $t = 2$: $B$ से आगे, उससे
$B$ की दूरी पर ($\vect{AC} = 2\vect{AB}$); $t = -1$: $A$ से होकर $B$ का परावर्तन।
\end{solution}

\begin{solution}{exo:b2:affine:3}
$M$ $(1,-1)$ के अनुदिश $\operatorname{Vect}(1,1)$ पर प्रक्षेप आव्यूह है ($M^2 = M$ जाँचिए)। $f$ का प्रतिबिंब:
$\{MX + C\} = C +
\operatorname{im} M$: अर्थात् $(1,1)$ से दिशित $(1,0)$ से होकर जाने वाली \omterm{def:b2:affine:subspace}{आनत} रेखा। अचल बिंदु: $X = MX + C$,
अर्थात् $(I - M)X = C$; परंतु $C =
(1, 0)^{\mathsf T}$ और $\operatorname{im}(I - M) =
\operatorname{Vect}(1,-1)$; क्या $(1,0)$ उसमें है? $(1, 0) =
\alpha(1,-1)$ $\alpha = 1$ तथा $0 = -1$ को
बाध्य कर देता है: नहीं। अतः कोई अचल बिंदु नहीं। और
\[
f(f(X)) = M(MX + C) + C = MX + MC + C = f(X) + MC,
\qquad MC = \tfrac12(1,1)^{\mathsf T} :
\]
$f\circ f$ $f$ के बाद प्रतिबिंब-रेखा के अनुदिश स्थानांतरण है --- अर्थात् $f$
एक ``सर्पी प्रक्षेप'' है: रेखा पर प्रक्षेप के साथ संयोजित सर्पण।
\end{solution}

\begin{solution}{exo:b2:affine:4}
$I = \operatorname{bar}(B, 2; C, 1)$ (क्योंकि $\vect{BI} =
\frac13\vect{BC}$ $I$ को $B$ के अधिक निकट रखता है: $B$ पर भार $2$, $C$ पर
$1$ --- जाँचिए: $\vect{BI} = \frac{1}{3}\vect{BC}$)। इसी प्रकार $J = \operatorname{bar}(C, 2; A, 1)$, $K =
\operatorname{bar}(A, 2; B, 1)$। तीनों भारित निकाय जोड़ने पर
$(I, 1; J, 1; K, 1)$ का \omterm{def:b2:affine:barycenter}{बैरिकेंद्र} (हर एक का कुल भार $3$ है, अतः $I$ के स्थान पर उसका
निकाय रखिए, इत्यादि)
\[
\operatorname{bar}\bigl(A, 1 + 2;\ B, 2 + 1;\ C, 1 + 2\bigr)
= \operatorname{bar}(A, 1; B, 1; C, 1) = G ,
\]
है, अर्थात् $ABC$ का गुरुत्व केंद्र: इस प्रकार त्रिभुज $IJK$ का गुरुत्व
केंद्र वही है --- और यह अनुमेय था, क्योंकि रचना $A, B, C$ को \omterm{def:b2:structures:generated}{चक्रीय} रूप से
बरतती है और गुरुत्व केंद्र भारों की \omterm{def:b2:structures:generated}{चक्रीय} सममिति का अद्वितीय अचल बिंदु है।
\end{solution}

\begin{solution}{exo:b2:affine:5}
$g(u) = f(O + u) - f(O)$ रखिए ($O$ पर सदिशीकृत $\R^n$ में काम करते हुए), $g(0) = 0$।

\emph{योज्यता:} $\frac{(O + u) + (O + v)}{2} = O + \frac{u +
v}{2}$, अतः मध्यबिंदु का सुरक्षित रहना $g\bigl(\frac{u+v}{2}\bigr)
= \frac{g(u) + g(v)}{2}$ देता है; और $v = 0$ के
साथ: $g(u/2) = g(u)/2$; दोनों मिलाकर $g(u + v) = 2g\bigl(\frac{u+v}{2}\bigr) = g(u) + g(v)$।

\emph{$\Q$-समघातता:} योज्यता $g(nu) = ng(u)$ देती है ($n \in
\N$, आगमन), फिर $g(-u) = -g(u)$ (जोड़िए), फिर
$g(\frac pq u) =
\frac pq g(u)$ ($q$ लगाइए, और मापन के एकैकीपन का उपयोग कीजिए)।

\emph{$\R$-समघातता:} $t \in \R$ के लिए परिमेय $t_n \to
t$ लीजिए: $g(t_nu) = t_ng(u)$, और $g$ की संततता
(जो $f$ से विरासत में मिली) सीमा तक ले जाती है: $g(tu) = tg(u)$। अतः $g$ रैखिक है और
$f =
f(O) + g$: यानी \omterm{def:b2:affine:subspace}{आनत}।
\end{solution}

\begin{solution}{exo:b2:affine:6}
रैखिक भाग $\varphi(x,y) = (y, x)$: विकर्ण $y = x$ में परावर्तन (लांबिक, \omterm{def:b2:linalg:det}{सारणिक} $-1$)। अचल बिंदु:
$(x, y) = (y
+ 1, x - 1)$ अकेले समीकरण $y = x - 1$ पर सिमट आता है (दोनों घटक तुल्य हैं): अतः रेखा $y = x - 1$ का
हर बिंदु अचल है। इस प्रकार $f$ उस रेखा को \omterm{def:b2:funcseq:def}{बिंदुवार} अचल रखता है: $f$ उसी
अक्ष में \emph{परावर्तन} है (अचल बिंदुओं की एक रेखा तथा परावर्तन रैखिक भाग
वाला समदूरिक प्रतिचित्रण)। और इसके अनुरूप $f \circ
f(x,y) = f(y+1, x-1) = (x - 1 + 1, y + 1 - 1) = (x, y)$: अर्थात् अंतर्वलन, जैसा
किसी परावर्तन को होना ही चाहिए।
\end{solution}

\begin{solution}{exo:b2:affine:7}
$i \geq 2$ सदिश $\vect{A_1A_i}$ ($n + 1$) विमा $n$ में परस्पर जुड़े हैं: अतः ऐसे $\mu_i$ हैं जो सब
शून्य नहीं और $\sum_{i\geq2}
\mu_i\vect{A_1A_i} = 0$; $\mu_1 = -\sum_{i \geq 2}\mu_i$ रखिए, तब $\sum_{i}\mu_i = 0$ और प्रत्येक $O$ के लिए $\sum_i \mu_i\,\vect{OA_i} = 0$, जहाँ सारे
$\mu_i$ शून्य नहीं। सूचकांकों को बाँटिए: $P = \{i : \mu_i >
0\}$, $N = \{i : \mu_i < 0\}$, और दोनों अरिक्त ($\mu_i$ का
योग शून्य है और वे सब शून्य नहीं)। $s = \sum_{i\in P}\mu_i =
-\sum_{i \in N}\mu_i > 0$ के साथ:
\[
\operatorname{bar}\bigl(A_i, \tfrac{\mu_i}{s}\bigr)_{i \in P}
= \operatorname{bar}\bigl(A_i, \tfrac{-\mu_i}{s}\bigr)_{i \in N},
\]
(दोनों पक्ष उसी संबंध से $\vect{OX} = \frac1s\sum_{i\in
P}\mu_i\vect{OA_i}$ वाला बिंदु $X$ देते हैं): अर्थात् दोनों
\omterm{def:b2:affine:convex}{उत्तल} कवचों का कोई साझा बिंदु, और सूचकांक-समूह असंयुक्त।
\end{solution}

\begin{solution}{exo:b2:affine:8}
\emph{प्रतिबिंब = अचल बिंदु:} $Y = f(X)$ के लिए $f(Y) = f(f(X)) =
f(X) = Y$: अर्थात् हर प्रतिबिंब-बिंदु
अचल है; और विलोमतः अचल बिंदु प्रतिबिंब हैं। अतः $\mathcal{F} = \operatorname{im} f =
\operatorname{Fix}(f)$ अरिक्त है, और वह \omterm{def:b2:affine:subspace}{आनत}
उपसमष्टि है (किसी \omterm{def:b2:affine:subspace}{आनत प्रतिचित्रण} का प्रतिबिंब) जिसकी दिशा $\operatorname{im}\vec f$ है।

\emph{प्रक्षेप संरचना:} $\vec f$ वर्गसम है ($\vec{f\circ
f} = \vec f^{\,2} = \vec f$), अतः $E = \operatorname{im}\vec f
\oplus \ker \vec f$ (\cref{ex:b2:reduction:projections})। किसी भी बिंदु
$X$ के लिए सदिश $\vect{f(X)\,X}$ लीजिए; $\vec f$ लगाने पर:
\[
\vec f\bigl(\vect{f(X)\,X}\bigr) = \vect{f(f(X))\,f(X)} = 0
\qquad (f \circ f = f),
\]
अतः $\vect{f(X)\,X} \in \ker\vec f$। इसलिए $X = f(X) +
\vect{f(X)X}$ $X$ को $\mathcal{F}$ के किसी बिंदु के रूप में दिखा देता है जो
$\ker\vec f$ के किसी सदिश से स्थानांतरित है: अर्थात् $f$ ठीक $\ker\vec f$ के अनुदिश $\mathcal{F}$
पर प्रक्षेप है। और विलोमतः ऐसे प्रक्षेप स्पष्टतः $f \circ f = f$ पूरा करते हैं।
\end{solution}

\begin{solution}{exo:b2:affine:9}
मान लीजिए $M = \operatorname{bar}(B, 2;\ C, 3)$, जिसका कुल भार $5$ है। साहचर्यता $G = \operatorname{bar}(A, 1;\ M, 5)$ देती है, अतः $\vect{AG} = \frac56\,\vect{AM}$:
अर्थात् $G$ खंड $\intcc AM$ पर $A$ से उसके पाँच-छठे भाग पर पड़ता है। और चूँकि $A \notin
(BC)$,
रेखा $(AG) = (AM)$ $(BC)$ से अकेले बिंदु $M$ पर मिलती है, जहाँ $\vect{BM} = \frac35\,\vect{BC}$। इसी प्रकार $N =
\operatorname{bar}(C, 3;\ A, 1)$
(कुल भार $4$, $\vect{CN}
= \frac14\,\vect{CA}$) के साथ साहचर्यता $G =
\operatorname{bar}(B, 2;\ N, 4)$ देती है: अतः रेखा $(BG)$ $(CA)$ से
$N$ पर मिलती है, और $\vect{BG} = \frac46\,\vect{BN} =
\frac23\,\vect{BN}$।
\end{solution}

\begin{solution}{exo:b2:affine:10}
किसी मूल बिंदु $O$ पर सदिशीकरण कीजिए और बिंदुओं को सदिश लिखिए: $\omega = \vect{O\Omega}$ के साथ
$h_{\Omega, \lambda}(x) = \omega + \lambda(x - \omega)$। संयोजन $g = h_{\Omega',
\mu} \circ h_{\Omega, \lambda}$ रैखिक भाग $\mu\lambda\,\mathrm{id}$ के साथ \omterm{def:b2:affine:subspace}{आनत} है। यदि $\lambda\mu \neq 1$: तो $1 \notin
\operatorname{Sp}(\lambda\mu\,\mathrm{id})$, अतः \cref{prop:b2:affine:fixedpoint}
अद्वितीय अचल बिंदु $\Omega''$ दे देता है और वहाँ सदिशीकरण करने पर $g =
\lambda\mu\,\mathrm{id}$: अर्थात्
समरूपण $h_{\Omega'',
\lambda\mu}$। और यदि $\lambda\mu = 1$, तो रैखिक भाग तत्समक है, अतः $g$ स्थानांतरण है;
और खोलने पर
\[
g(x) = \omega' + \mu\bigl(\omega + \lambda(x - \omega) -
\omega'\bigr) = x + (1 - \mu)\,\omega' + \mu(1 -
\lambda)\,\omega .
\]
$\lambda = \mu = -1$ (बिंदु परावर्तन) के लिए सदिश $2\omega' - 2\omega = 2\,\vect{\Omega\Omega'}$ है: अर्थात् पहले $\Omega$ में और फिर
$\Omega'$ में बिंदु परावर्तनों का संयोजन $2\,\vect{\Omega\Omega'}$ से स्थानांतरण है।
\end{solution}

\begin{solution}{exo:b2:affine:11}
$\vect{A'B} = \alpha\,\vect{A'C}$ ठीक $1\cdot
\vect{A'B} - \alpha\,\vect{A'C} = 0$ कहता है, अर्थात् $A' =
\operatorname{bar}(B, 1;\ C, -\alpha)$ (कुल भार $1 -
\alpha \neq 0$, क्योंकि $B \neq C$); इसी प्रकार
$B' =
\operatorname{bar}(C, 1;\ A, -\beta)$ और $C' =
\operatorname{bar}(A, 1;\ B, -\gamma)$।

\emph{संरेखता की कसौटी।} हर बिंदु $P$ को $(A, B, C)$ के सापेक्ष उसकी सामान्यीकृत
बैरिकेंद्रीय पंक्ति $p = (p_A, p_B, p_C)$, $p_A + p_B +
p_C = 1$ दीजिए। यदि तीन बिंदुओं $P_1, P_2,
P_3$ के लिए $(c_1, c_2, c_3) \neq 0$ सहित
$\sum_i c_i p_i = 0$ हो, तो प्रविष्टियाँ जोड़ने पर $\sum c_i = 0$ मिलता है, और $\sum_i c_i \vect{OP_i} = \sum_j \bigl(\sum_i
c_ip_{ij}\bigr)\vect{OV_j} = 0$: अर्थात् $P_i$ \omterm{def:b2:affine:subspace}{आनत}
रूप से परतंत्र, यानी संरेख हैं। विलोमतः कोई \omterm{def:b2:affine:subspace}{आनत} परतंत्रता $(t_i)$ $0$ तक
जुड़ने वाली प्रविष्टियों तथा $\sum_j w_j\vect{OV_j} = 0$ के साथ $w = \sum t_ip_i$ देती है; और $A$ से खोलने पर
$w_B
\vect{AB} + w_C\vect{AC} = 0$, अतः $(A, B, C)$ की \omterm{def:b2:affine:subspace}{आनत} स्वतंत्रता से $w = 0$: यानी पंक्तियाँ रैखिकतः परतंत्र
हैं। इस प्रकार संरेखता किसी $3 \times 3$ \omterm{def:b2:linalg:det}{सारणिक} के लुप्त होने पर सिमट आती है, और
पंक्तियों को अशून्य गुणकों $1 - \alpha$, $1 -
\beta$, $1 - \gamma$ से मापित करने से कुछ नहीं
बदलता:
\[
\det\begin{pmatrix} 0 & 1 & -\alpha\\ -\beta & 0 & 1\\ 1 &
-\gamma & 0\end{pmatrix}
= 1 - \alpha\beta\gamma .
\]
अतः $A', B', C'$ संरेख हैं तभी जब $\alpha\beta\gamma = 1$: यही मेनेलौस की प्रमेय है।
\end{solution}

\begin{solution}{exo:b2:affine:12}
मान लीजिए $\Delta = \{\lambda \in \R^{n+1} : \lambda_i \geq 0,\
\sum\lambda_i = 1\}$: यह $\R^{n+1}$ में संवृत और परिबद्ध है, अतः \omterm{def:b2:metric:compact}{संहत}, और $K^{n+1}$ परिमित
गुणन के रूप में \omterm{def:b2:metric:compact}{संहत} है। प्रतिचित्रण
\[
\Phi \colon \Delta \times K^{n+1} \to \R^n, \qquad
\Phi(\lambda, x_0, \dots, x_n) = \sum_{i=0}^n \lambda_i x_i
\]
\omterm{def:b2:metric:continuity}{संतत} है, और \cref{thm:b2:affine:caratheodory} ठीक यही कहता है कि $\operatorname{conv}(K) = \Phi(\Delta \times
K^{n+1})$: अर्थात् किसी \omterm{def:b2:metric:compact}{संहत} समुच्चय का \omterm{def:b2:metric:continuity}{संतत}
प्रतिबिंब (\cref{thm:b2:metric:compactprops}), इसलिए \omterm{def:b2:metric:compact}{संहत}।

किसी संवृत समुच्चय के लिए: $S = (\R \times \{0\}) \cup \{(0,
1)\}$ लीजिए, जो $\R^2$ में संवृत है। $(0,1)$ पर भार $t$
और अक्ष-बिंदुओं पर $1 - t$ डालने वाले किसी उत्तल संचय का दूसरा निर्देशांक $t$
होता है, अतः
\[
\operatorname{conv}(S) = \bigl(\R \times \intco01\bigr) \cup
\{(0,1)\}
\]
($0 \leq t < 1$ के लिए $(x, t) = t\,(0,1) + (1-t)\,\bigl(\tfrac
x{1-t}, 0\bigr)$)। बिंदु $(1, 1) = \lim_{t \to 1}(1, t)$ संलग्न है पर कवच में नहीं: अतः संवृत नहीं।
\end{solution}

\begin{solution}{pb:b2:affine:1}
\textbf{1.} $A_j$ पर आधार बदलिए: $i \neq j$ के लिए $\vect{A_jA_i} =
\vect{A_0A_i} - \vect{A_0A_j}$। यदि $\sum_{i \neq j}
c_i\vect{A_jA_i} = 0$, तो खोलने पर $\sum_{i \notin \{0,
j\}} c_i\,\vect{A_0A_i} - \bigl(\sum_{i\neq j}
c_i\bigr)\vect{A_0A_j} = 0$;
और $\vect{A_0A_i}$ की स्वतंत्रता $i \notin \{0, j\}$ के लिए $c_i = 0$ को बाध्य कर देती है, फिर $c_0 = 0$: अर्थात्
$A_j$ पर स्वतंत्रता। तुल्यता के लिए: $1$ तक जुड़ने वाले तथा एक ही
\omterm{def:b2:affine:barycenter}{बैरिकेंद्र} वाले दो भार-कुल $(\lambda_i)$, $(\lambda_i')$ $\nu = \lambda -
\lambda'$ के साथ $\sum\nu_i = 0$ और (मूल बिंदु $A_0$
से) $\sum_{i \geq
1}\nu_i\,\vect{A_0A_i} = 0$ देते हैं, अतः स्वतंत्रता के अंतर्गत $\nu = 0$। विलोमतः $\mu_0 = -\sum_{i\geq1}
\mu_i$ से पूरा
किया गया कोई अतुच्छ संबंध $\sum_{i\geq1}\mu_i
\vect{A_0A_i} = 0$ \omterm{def:b2:affine:barycenter}{बैरिकेंद्र} हिलाए बिना किसी भी भार-कुल में
$t(\mu_i)$ जोड़ने देता है: अर्थात् अनद्वितीयता।

\textbf{2.} $(\vect{AB}, \vect{AC})$ $\R^2$ का आधार है: $\vect{AM} = \beta\,\vect{AB} + \gamma\,\vect{AC}$ (अद्वितीय) लिखिए और $\alpha = 1 - \beta - \gamma$ रखिए; और मूल
बिंदु $A$ पर बैरिकेंद्र-शर्त ठीक $\vect{AM} =
\beta\,\vect{AB} + \gamma\,\vect{AC}$ पढ़ी जाती है। अद्वितीयता प्रश्न 1 है।

\textbf{3.} $\alpha\vect{MA} + \beta\vect{MB} +
\gamma\vect{MC} = 0$ और शाल से $\vect{MA} =
-\beta\,\vect{AB} - \gamma\,\vect{AC}$। $D =
\det(\vect{AB}, \vect{AC})$ के साथ:
\begin{align*}
\det(\vect{MB}, \vect{MC})
&= \det(\vect{MA} + \vect{AB},\ \vect{MA} + \vect{AC})\\
&= \det(\vect{MA}, \vect{AC}) + \det(\vect{AB}, \vect{MA}) +
D = -\beta D - \gamma D + D = \alpha D,
\end{align*}
जिसमें द्विरैखिकता तथा $\det(\vect{MA}, \vect{AC}) = -\beta D$, $\det(\vect{AB}, \vect{MA}) = -\gamma D$ का उपयोग हुआ। शेष दोनों सूत्र भूमिकाओं को
\omterm{def:b2:structures:generated}{चक्रीय} रूप से बदलकर उसी परिकलन से निकलते हैं।

\textbf{4.} परिभाषा से $\operatorname{conv}\{A, B, C\}$ अऋणात्मक भारों वाले बैरिकेंद्रों का समुच्चय है;
और भारों को $1$ तक सामान्यीकृत करके अद्वितीयता (प्रश्न 2) का सहारा लेने
पर $M \in \operatorname{conv}\{A,B,C\}$ तभी जब $\alpha, \beta,
\gamma \geq 0$। हर निर्देशांक $M$ का \omterm{def:b2:affine:subspace}{आनत} फलन है (प्रश्न 3: $M$ में \omterm{def:b2:affine:subspace}{आनत}
एक स्तंभ वाला $2\times2$ \omterm{def:b2:linalg:det}{सारणिक}), अतः हर \omterm{def:b2:metric:topology}{विवृत} चिह्न-शर्त कोई \omterm{def:b2:metric:topology}{विवृत} अर्धसमतल
परिभाषित करती है। प्रतिरूप $(-,-,-)$ $\alpha + \beta
+ \gamma = 1$ के विरुद्ध है; और शेष सातों प्रतिरूप
साकार होते हैं: चिह्नों का पालन करने वाले किसी त्रिक को, जिसमें कम से कम
एक $+$ प्रविष्टि हो, इस प्रकार मापित कीजिए कि (धनात्मक) योग $1$ हो जाए ---
जैसे $(-1, 1, 1)$, $(3, -1, -1)$, $(\frac13, \frac13, \frac13)$, और उनके क्रमचय।

\textbf{5.} कोई \omterm{def:b2:affine:subspace}{आनत प्रतिचित्रण} \omterm{def:b2:affine:barycenter}{बैरिकेंद्र} सुरक्षित रखता है (\cref{def:b2:affine:subspace}), अतः
$u(M) = \alpha u(A) +
\beta u(B) + \gamma u(C)$। $(a, b) \neq (0,0)$ के साथ $u(x, y) = ax + by + c$ लिखने पर: $\{u = c'\}$ कोई रेखा है, और हर रेखा $ax + by = c'$ ऐसा ही स्तर
समुच्चय है। और यदि $u(M), u(N) \geq
c$ तथा $t \in \intcc01$, तो $u\bigl(\operatorname{bar}(M,
1-t; N, t)\bigr) = (1-t)u(M) + tu(N) \geq c$: अर्थात् अर्धसमतल \omterm{def:b2:affine:convex}{उत्तल} हैं।

\textbf{6.} शर्तें $\sum\mu_i = 0$, $3\mu_2 + \mu_4 =
0$, $3\mu_3 + \mu_4 = 0$ ($\mu_4 = 3$ लेकर) परतंत्रता $(\mu_1, \mu_2, \mu_3, \mu_4) = (-1, -1, -1, 3)$ देती हैं। चिह्न
$\{A_4\} \mid \{A_1, A_2, A_3\}$ के रूप में बँटते हैं, और हर पक्ष को $3$ से सामान्यीकृत करने पर:
\[
A_4 = \operatorname{bar}\bigl(A_1, \tfrac13;\ A_2, \tfrac13;\
A_3, \tfrac13\bigr) = (1,1),
\]
अर्थात् त्रिभुज का गुरुत्व केंद्र: रादों बिंदु स्वयं $A_4$ है, जो वास्तव में
त्रिभुज $A_1A_2A_3$ के भीतर पड़ता है।

\textbf{7.} रैखिक प्रतिचित्रण $\Phi \colon \R^{n+2} \to \R
\times \R^n$, $\mu \mapsto (\sum\mu_i,\
\sum\mu_i\vect{OA_i})$, की कोटि $\leq n + 1$ है, अतः $\dim\ker
\Phi \geq 1$। यदि दो
स्वतंत्र परतंत्रताएँ $\mu, \mu'$ होतीं, तो कोई उपयुक्त संचय $\nu = \mu'_{n+2}\mu -
\mu_{n+2}\mu'$ (अथवा $\mu$ स्वयं,
यदि दोनों के अंतिम गुणांक लुप्त हों) $\nu_{n+2}
= 0$ वाली \emph{अशून्य} परतंत्रता
होती; और $A_1, \dots, A_{n+1}$ तक प्रतिबंधित करके $A_1$ पर आधार बदलने पर कोई $\nu_i \neq 0$ $i \geq 2$ के साथ
मिलता (अकेला अशून्य भार शून्य तक नहीं जुड़ सकता), जिससे अतुच्छ संबंध $\sum_{i\geq2}\nu_i\vect{A_1A_i} = 0$
मिलता: अर्थात् $n+1$ बिंदु \omterm{def:b2:affine:subspace}{आनत} रूप से परतंत्र होते, जो व्यापक स्थिति के
विरुद्ध है। अतः $\dim\ker
\Phi = 1$। और वही प्रतिबंधन-तर्क दिखाता है कि किसी अशून्य
परतंत्रता का कोई गुणांक लुप्त नहीं होता।

\textbf{8.} मान लीजिए $\mu \neq 0$ कोई परतंत्रता है, $P = \{i :
\mu_i > 0\}$ और $N = \{i : \mu_i < 0\}$: दोनों अरिक्त ($\sum\mu_i = 0$,
$\mu \neq 0$) और \omterm{def:b2:metric:complete}{पूर्ण} (कोई गुणांक शून्य नहीं)। रादों की रचना (\cref{exo:b2:affine:7}) ठीक इसी विभाजन
से साझा कवच-बिंदु उत्पन्न करती है। और चूँकि परतंत्रता किसी अशून्य अदिश तक
अद्वितीय है (प्रश्न 7), अक्रमित युग्म $\{P, N\}$ --- अतः रादों विभाजन --- अद्वितीय
है।

\textbf{9.} खंड अरिक्त हैं, अतः प्रकार $(1,3)$ या $(2,2)$ है। प्रकार $(1,3)$, खंड $\{j\}$:
रादों बिंदु $\operatorname{conv}\{A_j\} = \{A_j\}$ में पड़ता है, अतः शेष तीन का $A_j \in
\operatorname{conv}$; और वह किसी भुजा पर नहीं
पड़ सकता (तब तीन बिंदु संरेख हो जाते, जो व्यापक स्थिति के विरुद्ध है), अतः
$A_j$ त्रिभुज के भीतर है। प्रकार $(2,2)$, खंड $\{i,j\} \mid \{k,l\}$: रादों बिंदु $z$ $\intcc{A_i}{A_j} \cap \intcc{A_k}{A_l}$ पर पड़ता
है, और $z$ कोई अंत्यबिंदु नहीं (तब तीन बिंदु संरेख हो जाते): अतः दोनों
खंड किसी अभ्यंतरीय बिंदु पर कटते हैं। इसके अतिरिक्त कोई भी बिंदु शेष के
कवच में नहीं है: ऐसा समावेश $\lambda_i \geq 0$ के साथ $A_l =
\operatorname{bar}(A_i, \lambda_i)_{i \neq l}$ $(+,+,+,-)$ चिह्न-प्रतिरूप वाली \omterm{def:b2:affine:subspace}{आनत}
परतंत्रता होती, जो अद्वितीयता (प्रश्न 8) से विभाजन को $(1,3)$ बना देती। अतः
$(2,2)$ वाली स्थिति में चारों बिंदु \omterm{def:b2:affine:convex}{उत्तल} स्थिति में हैं और कटते हुए खंड
विकर्ण हैं।

\textbf{10.} समीकरण $\mu_2 + \mu_3 = 0$, $\mu_3 +
\mu_4 = 0$, $\sum\mu_i = 0$ परतंत्रता $(1, -1, 1,
-1)$ देते हैं: विभाजन $\{(0,0), (1,1)\} \mid \{(1,0), (0,1)\}$, और
\[
\operatorname{bar}\bigl((0,0), \tfrac12;\ (1,1),
\tfrac12\bigr) = \bigl(\tfrac12, \tfrac12\bigr) =
\operatorname{bar}\bigl((1,0), \tfrac12;\ (0,1),
\tfrac12\bigr) :
\]
अर्थात् रादों बिंदु वर्ग का केंद्र है, जहाँ दोनों विकर्ण कटते हैं ---
प्रकार $(2,2)$, जैसा चित्र बताता है।

\textbf{11.} $x_1, \dots, x_4$ पर लगाई गई रादों खंड $I \mid J$ तथा कोई बिंदु $z \in \operatorname{conv}\{x_i : i
\in I\} \cap \operatorname{conv}\{x_j : j \in J\}$ देती है। $k
\in \{1, \dots, 4\}$
स्थिर कीजिए, मान लीजिए $k \in I$। हर $j \in J$ $j \neq k$ पूरा करता है, अतः $x_j \in
\bigcap_{l \neq j}C_l$ के चुनाव से
$x_j \in C_k$; और चूँकि $C_k$ \omterm{def:b2:affine:convex}{उत्तल} है, $z \in
\operatorname{conv}\{x_j : j \in J\} \subseteq C_k$। और $k$ स्वेच्छ था, अतः $z \in C_1 \cap C_2 \cap C_3 \cap C_4$।

\textbf{12.} $m$ पर आगमन। $m = 3$ के लिए परिकल्पना ही निष्कर्ष है; और $m = 4$
प्रश्न 11 है। मान लीजिए $m \geq 4$, कथन $m$ समुच्चयों के लिए मान लीजिए, और
त्रिक-प्रतिच्छेदन गुणधर्म वाले $C_1, \dots,
C_{m+1}$ लीजिए। $C_m' =
C_m \cap C_{m+1}$ रखिए, जो \omterm{def:b2:affine:convex}{उत्तल} है। कुल $C_1, \dots, C_{m-1},
C_m'$
के $m$ सदस्य हैं; $C_m'$ से बचने वाला कोई त्रिक परिकल्पना से प्रतिच्छेद करता
है, और त्रिक $\{C_i, C_j, C_m'\}$ का प्रतिच्छेद $C_i \cap C_j \cap C_m \cap C_{m+1}$ है, जो $C_i, C_j, C_m, C_{m+1}$ पर लगाए गए प्रश्न 11 से
अरिक्त है (इनमें से कोई भी तीन परिकल्पना से मिलते हैं)। अब आगमन परिकल्पना
नए कुल का, अर्थात् सभी $m+1$ समुच्चयों का, साझा बिंदु दे देती है।

\textbf{13.} (क) किसी अनपभ्रष्ट त्रिभुज की संवृत भुजाएँ $\intcc PQ$, $\intcc QR$, $\intcc RP$:
कोई भी दो कोई शीर्ष साझा करती हैं, पर तीनों का साझा बिंदु $\intcc
PQ \cap \intcc RP = \{P\}$ तथा $\intcc QR$ में
होता, जो $P$ को अपवर्जित करता है। (ख) समुच्चयों $S_i = \{x_1, \dots, x_4\}
\setminus \{x_i\}$ में से कोई भी तीन
चारों बिंदुओं में से तीन छोड़ देते हैं, जिससे ठीक एक साझा बिंदु बचता है;
पर कुल प्रतिच्छेद हर बिंदु छोड़ देता है। $S_i$ परिमित हैं, \omterm{def:b2:affine:convex}{उत्तल} नहीं:
अर्थात् उत्तलता अनिवार्य है। (ग) परिमित $H_k = \intco{k}{+\infty} \times
\R$ का प्रतिच्छेद $\intco{k_{\max}}{+\infty} \times \R \neq
\emptyset$ है, फिर भी
किसी बिंदु के लिए सभी $k \in \N$ हेतु $x \geq k$ नहीं: अनंत कुलों के लिए संहतता
अनिवार्य है।

\textbf{14.} मान लीजिए $\bigcap_{i \in I}K_i = \emptyset$ और $i_0$ स्थिर कीजिए। हर $x \in K_{i_0}$ किसी न किसी $K_i$ से
चूक जाता है, अतः $K_{i_0} \subseteq \bigcup_{i \in I}(\R^2 \setminus K_i)$, जो \omterm{def:b2:metric:topology}{विवृत} समुच्चयों से आच्छादन है ($K_i$ \omterm{def:b2:metric:compact}{संहत} है, अतः
संवृत)। बोरेल--लेबेग (\cref{thm:b2:metric:borellebesgue}) से परिमित ही पर्याप्त हैं: $K_{i_0} \cap K_{i_1} \cap \dots \cap K_{i_N} =
\emptyset$। परंतु \omterm{def:b2:affine:convex}{उत्तल}
समुच्चयों के इस परिमित कुल के किन्हीं तीन सदस्य प्रतिच्छेद करते हैं, अतः
प्रश्न 12 प्रतिच्छेद को अरिक्त बना देता है: विरोधाभास।

\textbf{15.} $p \in S$ के लिए $D_p = \overline D(p, r)$ रखिए: ये \omterm{def:b2:metric:compact}{संहत} \omterm{def:b2:affine:convex}{उत्तल समुच्चय} हैं। $p, q, s \in S$ के लिए
परिकल्पना $p, q, s$ को समाहित करने वाली कोई संवृत चक्रिका $\overline D(z, r)$ देती है; तब $\norm{\vect{zp}}, \norm{\vect{zq}}, \norm{\vect{zs}}
\leq r$,
अर्थात् $z \in D_p \cap D_q \cap D_s$। हेली (प्रश्न 12; कुल परिमित है) से कोई $c \in
\bigcap_{p \in S}D_p$ है: अतः हर $p \in S$
$\norm{\vect{cp}} \leq r$ पूरा करता है, यानी $S \subseteq \overline D(c,
r)$।

\textbf{16.} मान लीजिए $c = x_{\lceil n/2\rceil}$। $c$ को समाहित करने वाली कोई संवृत अर्धरेखा
$t \geq
c$ के साथ $\intoc{-\infty}{t}$ है अथवा $t \leq c$ के साथ $\intco{t}{+\infty}$। पहली $x_1, \dots, x_{\lceil n/2\rceil}$ को समाहित करती है: अर्थात्
कम से कम $\lceil n/2\rceil \geq n/2$ बिंदु। और दूसरी $x_{\lceil n/2\rceil}, \dots, x_n$ को: ठीक $n - \lceil
n/2\rceil + 1 = \floor{n/2} + 1 > n/2$ बिंदु।

\textbf{17.} $\abs{A \cap B} = \abs A + \abs B - \abs{A \cup
B} \geq \abs A + \abs B - n > \tfrac{4n}3 - n = \tfrac n3$, अतः
\[
\abs{A \cap B \cap C} \geq \abs{A \cap B} + \abs C - n >
\tfrac n3 + \tfrac{2n}3 - n = 0 .
\]

\textbf{18.} ध्यान दीजिए $m > 2n/3$। गणनसंख्या $m$ वाले $T_1, T_2, T_3 \subseteq S$ के लिए प्रश्न 17 कोई
बिंदु $x \in T_1
\cap T_2 \cap T_3$ देता है; तब हर $i$ के लिए $x \in \operatorname{conv}(T_i)$: अर्थात् $\mathcal F$ के किन्हीं तीन सदस्य
मिलते हैं। कुल परिमित है ($S$ के परिमित उपसमुच्चय) और \omterm{def:b2:affine:convex}{उत्तल} समुच्चयों से
बना है, अतः हेली (प्रश्न 12) $c \in
\bigcap_{\abs T = m}\operatorname{conv}(T)$ दे देती है।

\textbf{19.} मान लीजिए किसी संवृत अर्धसमतल $H \ni c$ में $S$ के $n/3$ से कम बिंदु
हैं। उसका पूरक $U$ \omterm{def:b2:metric:topology}{विवृत} अर्धसमतल है, \omterm{def:b2:affine:convex}{उत्तल}, जहाँ $\abs{S \cap U} > 2n/3$, अतः $\abs{S \cap U} \geq \floor{2n/3} + 1 = m$; और $\abs T = m$
वाला $T
\subseteq S \cap U$ चुनिए। तब $U$ की उत्तलता से $\operatorname{conv}(T) \subseteq U$, अतः $c \in \operatorname{conv}(T) \subseteq U$: जो $c \in H$ के विरुद्ध है।
इसलिए $c$ को समाहित करने वाला हर संवृत अर्धसमतल कम से कम $n/3$ बिंदु
समाहित करता है: अर्थात् $c$ केंद्रबिंदु है।

\textbf{20.} भुजा $L$ वाला समबाहु त्रिभुज और $\varepsilon = L/100$ लीजिए। मान लीजिए $c$ कोई
बिंदु है; हम $c$ को समाहित करने वाला ऐसा संवृत अर्धसमतल दिखाते हैं जिसमें
अधिकतम $k$ बिंदु हों।

\emph{स्थिति 1: $c$ किसी शीर्ष, मान लीजिए $B$, से $L/10$ के भीतर है।} $c$ से
$A$ तथा $C$ की दिशाएँ दिशाओं $B \to A$, $B \to C$ से अधिकतम
$\arcsin\bigl(\tfrac{L/10}{9L/10}\bigr) = \arcsin\tfrac19$
हटती हैं, अतः उनके बीच का कोण $\leq \tfrac\pi3 + 2\arcsin\tfrac19 <
\tfrac{2\pi}3$ है। \emph{स्थिति 2: $c$ सारे शीर्षों से
$> L/10$ दूरी पर है।} यदि $c$ त्रिभुज में हो, तो $c$ से शीर्षों की दिशाओं
$u_A, u_B, u_C$ के बीच के तीनों कोणीय अंतराल $2\pi$ तक जुड़ते हैं, अतः कोई अंतराल $\leq 2\pi/3$
है; और यदि $c$ बाहर हो, तो तीनों दिशाएँ दिशाओं के किसी \omterm{def:b2:metric:topology}{विवृत} अर्धसमतल में
पड़ती हैं और उनमें से दो का कोण $<
\pi/2$ होता है। हर स्थिति में दो दिशाएँ, मान
लीजिए $X$ और $Y$ की ओर, कोण $\leq 2\pi/3$ बनाती हैं; मान लीजिए $w$ उनका इकाई
समद्विभाजक है, अतः $\langle u_X, w\rangle, \langle u_Y, w\rangle \geq
\cos\tfrac\pi3 = \tfrac12$। $X$ के चारों ओर वाली चक्रिका के किसी भी बिंदु
$b$ के लिए:
\[
\langle \vect{cb}, w\rangle \geq
\tfrac12\norm{\vect{cX}} - \varepsilon > 0,
\]
क्योंकि $\norm{\vect{cX}} \geq L/10 > 2\varepsilon$ (और $Y$ के लिए भी वैसा ही): अतः \omterm{def:b2:metric:topology}{विवृत} अर्धसमतल $\{\langle \vect{cx},
w\rangle > 0\}$ दोनों
गुच्छ निगल लेता है। और उसका संवृत पूरक $c$ तथा अधिकतम तीसरे गुच्छ के $k$
बिंदु समाहित करता है। इस प्रकार समतल का कोई बिंदु $n/3$ को नहीं हरा पाता:
और प्रश्न 19 के साथ केंद्रबिंदु अचर ठीक $1/3$ है।

\textbf{21.} कोणों को क्रमित कीजिए; सबसे बड़ा $\theta$ $\theta \geq \pi/3$ पूरा करता है (तीनों
का योग $\pi$ है)। यदि $\theta \geq \pi/2$, मान लीजिए $R$ पर, तो $M$ को सामने वाली भुजा $\intcc PQ$
का मध्यबिंदु लीजिए। माध्यिका सूत्र ($\vect{RM}
= \tfrac12(\vect{RP} + \vect{RQ})$, खोलिए और कोसाइन नियम से $\langle\vect{RP}, \vect{RQ}\rangle$ हटा
दीजिए) देता है
\[
\norm{\vect{RM}}^2 = \tfrac12\norm{\vect{RP}}^2 +
\tfrac12\norm{\vect{RQ}}^2 - \tfrac14\norm{\vect{PQ}}^2 \leq
\tfrac12\norm{\vect{PQ}}^2 - \tfrac14\norm{\vect{PQ}}^2 =
\tfrac14\norm{\vect{PQ}}^2,
\]
जिसमें $\norm{\vect{PQ}}^2 = \norm{\vect{RP}}^2 +
\norm{\vect{RQ}}^2 - 2\langle\vect{RP}, \vect{RQ}\rangle \geq
\norm{\vect{RP}}^2 + \norm{\vect{RQ}}^2$ का उपयोग हुआ (अंतर्गुणन $\leq 0$ है)। अतः व्यास $\intcc PQ$ वाली, यानी
त्रिज्या $\leq \tfrac12 < \tfrac1{\sqrt3}$ वाली चक्रिका तीनों बिंदु समाहित कर लेती है (और अपभ्रष्ट
संरेख त्रिक $\theta = \pi$ के अंतर्गत आ जाते हैं)। और यदि $\theta < \pi/2$, तो त्रिभुज न्यूनकोण
है; ज्या-नियम से परित्रिज्या $R_c = a/(2\sin\theta)$ है, जहाँ $a \leq 1$ $\theta$ के सामने की भुजा है,
और $\theta \in
\intco{\pi/3}{\pi/2}$ $\sin\theta \geq \sqrt3/2$ देता है, अतः $R_c \leq 1/\sqrt3$: अर्थात् परिवृत्त काम कर देता है।

\textbf{22.} $p \in S$ के लिए $K_p = \overline D(p,
1/\sqrt3)$ रखिए: यह \omterm{def:b2:metric:compact}{संहत} \omterm{def:b2:affine:convex}{उत्तल} है। $S$ के किन्हीं तीन
बिंदु $p, q, s$ जोड़े-जोड़े में $\leq 1$ दूरी पर हैं, अतः प्रश्न 21 उन्हें समाहित
करने वाली त्रिज्या $1/\sqrt3$ की कोई चक्रिका देता है, जिसका केंद्र $K_p \cap K_q \cap K_s$ में पड़ता
है। \omterm{def:b2:metric:compact}{संहत} हेली (प्रश्न 14, जहाँ स्वेच्छ कुल अनुमत हैं) से कोई $c \in \bigcap_{p\in
S}K_p$ है: अतः
हर $p \in S$ $c$ से $1/\sqrt3$ के भीतर है, अर्थात् $S
\subseteq \overline D(c, 1/\sqrt3)$। यही समतल में युंग की प्रमेय
है।

\textbf{23.} गुरुत्व केंद्र $G$ के साथ $\sum_i\vect{GA_i} = 0$, अतः
\[
\sum_i\norm{\vect{OA_i}}^2 = \sum_i\norm{\vect{OG} +
\vect{GA_i}}^2 = 3\norm{\vect{OG}}^2 + 2\Bigl\langle
\vect{OG}, \sum_i\vect{GA_i}\Bigr\rangle +
\sum_i\norm{\vect{GA_i}}^2,
\]
और बीच वाला पद लुप्त हो जाता है: यही लाइब्नित्स सर्वसमिका है। भुजा $1$
वाले समबाहु त्रिभुज के लिए $\norm{\vect{GA_i}} =
1/\sqrt3$ (ऊँचाई $\sqrt3/2$ का दो-तिहाई), अतः $\sum_i\norm{\vect{GA_i}}^2 = 1$। और यदि
$\overline D(O, r)$ शीर्षों को समाहित करे, तो $3r^2 \geq
\sum_i\norm{\vect{OA_i}}^2 = 3\norm{\vect{OG}}^2 + 1 \geq 1$: $r \geq 1/\sqrt3$, जहाँ समता $O = G$ तथा तीनों दूरियों
के $r$ के बराबर होने को बाध्य कर देती है --- अर्थात् परिवृत्त। इस प्रकार
युंग का अचर $1/\sqrt3$ तीखा है।

\textbf{24.} \emph{$\R^n$ में हेली: यदि $C_1, \dots, C_m$ ($m \geq n + 1$) $\R^n$ के \omterm{def:b2:affine:convex}{उत्तल} उपसमुच्चय हों
और उनमें से कोई भी $n + 1$ प्रतिच्छेद करते हों, तो वे सब प्रतिच्छेद करते हैं।}
आधार स्थिति $m = n +
2$: $x_i \in \bigcap_{j\neq i}C_j$ चुनिए; रादों की प्रमेयिका (\cref{exo:b2:affine:7}) $x_1, \dots, x_{n+2}$ को खंड $I \mid J$ में
बाँट देती है जिनका साझा कवच-बिंदु $z$ है, और प्रत्येक $k$ के लिए जिस खंड
में $k$ नहीं है वह $C_k$ के बिंदुओं से बना है, अतः उत्तलता से $z \in C_k$, ठीक
जैसा प्रश्न 11 में। $m \geq n + 2$ के लिए आगमन-चरण: $C_m,
C_{m+1}$ के स्थान पर $C_m \cap C_{m+1}$ रखिए; नए
कुल का कोई ऐसा $(n+1)$-क्रम जिसमें प्रतिच्छेदित सदस्य हो, पुराने समुच्चयों के
$n + 2$ के बराबर है और आधार स्थिति उसे निपटा देती है, तथा शेष क्रम परिकल्पना
से ढँक जाते हैं। अब आगमन परिकल्पना से निष्कर्ष निकालिए।

\textbf{25.} रैखिक बीजगणित एक ही बार प्रवेश करता है: युग्मों (कुल भार,
भारित स्थिति) की $(n+1)$-विमीय समष्टि में $n + 2$ सदिश परतंत्र होने ही चाहिए ---
और वही \omterm{def:b2:affine:subspace}{आनत} परतंत्रता है। उसके गुणांकों के चिह्न बिंदुओं को दोनों रादों
खंडों में बाँट देते हैं और एक रैखिक संबंध को दो अऋणात्मक बैरिकेंद्रों की
समता में बदल देते हैं। उत्तलता ठीक दो बार काम आती है: हेली वाले पग में
($C_k$ के बिंदुओं का कवच $C_k$ में ही रहता है) और अनुप्रयोगों में (अर्धसमतल
और चक्रिकाएँ \omterm{def:b2:affine:convex}{उत्तल} हैं)। समतल की विमा केवल रादों को दिए गए बिंदुओं की
संख्या $4 = 2 + 2$ के द्वारा आई, अर्थात् हेली की परिकल्पना वाले ``$3 = 2 + 1$'' के द्वारा;
और शेष सब विमा-निरपेक्ष था, जैसा प्रश्न 24 पुष्ट करता है। $\R^n$ में अचर ये
बन जाते हैं: हेली संख्या $n + 1$; केंद्रबिंदु अचर $\frac1{n+1}$ (हर परिमित समुच्चय का
कोई ऐसा बिंदु होता है कि उससे होकर जाने वाला हर संवृत अर्ध-समष्टि उसका $\geq
\frac1{n+1}$
भाग समाहित करे); और व्यास $1$ वाले समुच्चयों के लिए युंग त्रिज्या $\sqrt{\frac{n}{2(n+1)}}$ ---
जो $n = 2$ होने पर $1/\sqrt3$ के बराबर है।
\end{solution}
